\documentclass[../../main.tex]{subfiles}
\graphicspath{{\subfix{../../image/}}}

\begin{document}

\section{习题}\label{section:ch05-习题}

在这些习题中，$U$ 总表示 $\mathbb{R}^{n}$ 的一个开子集，边界 $\partial U$ 光滑。如通常那样，除非另作说明，所有给定函数均假定光滑。

\begin{enumerate}
\item 设 $0<\beta<\gamma\leqslant 1$。证明插值不等式
\begin{align*}
\|u\|_{C^{0,\gamma}(U)}\leqslant\|u\|_{C^{0,\beta}(U)}^{\frac{1-\gamma}{1-\beta}}\|u\|_{C^{0,1}(U)}^{\frac{\gamma-\beta}{1-\beta}}.
\end{align*}

\item 记 $U$ 为开正方形 $\{x\in\mathbb{R}^{2}\mid |x_{1}|<1,\ |x_{2}|<1\}$。定义
\begin{align*}
u(x)=
\begin{cases}
1-x_{1}, & x_{1}>0,\ |x_{2}|<x_{1},\\
1+x_{1}, & x_{1}<0,\ |x_{2}|<-x_{1},\\
1-x_{2}, & x_{2}>0,\ |x_{1}|<x_{2},\\
1+x_{2}, & x_{2}<0,\ |x_{1}|<-x_{2}.
\end{cases}
\end{align*}
对哪些 $1\leqslant p\leqslant\infty$，$u$ 属于 $W^{1,p}(U)$？

\item 设 $U,V$ 为开集，$V\subset \subset U$。证明存在光滑函数 $\zeta$，使 $\zeta\equiv 1$ 于 $V$，$\zeta=0$ 于 $\partial U$ 附近。（提示：取 $V\subset \subset W\subset \subset U$，并对 $\chi_{W}$ 磨光。）

\item 设 $U$ 有界，且存在光滑向量场 $\boldsymbol{\alpha}$，使沿 $\partial U$ 有 $\boldsymbol{\alpha}\cdot\boldsymbol{\nu}\geqslant 1$，其中 $\boldsymbol{\nu}$ 如常表示单位外法向。设 $1\leqslant p<\infty$。对 $\int_{\partial U}|u|^{p}\boldsymbol{\alpha}\cdot\boldsymbol{\nu}\,\mathrm{d}S$ 应用 Gauss--Green 定理，导出迹不等式
\begin{align*}
\int_{\partial U}|u|^{p}\,\mathrm{d}S\leqslant C\int_{U}|Du|^{p}+|u|^{p}\,\mathrm{d}x
\end{align*}
的一个新证明，对一切 $u\in C^{1}(\overline{U})$ 成立。

\item 设 $U$ 有界，边界为 $C^{1}$。证明 $L^{p}(U)$（$1\leqslant p<\infty$）中一个“典型”函数 $u$ 在 $\partial U$ 上没有迹。更确切地说，证明不存在有界线性算子
\begin{align*}
T:L^{p}(U)\to L^{p}(\partial U)
\end{align*}
使每当 $u\in C(\overline{U})\cap L^{p}(U)$ 时 $Tu=u|_{\partial U}$。

\item 分部积分以证明插值不等式
\begin{align*}
\|Du\|_{L^{2}}\leqslant C\|u\|_{L^{2}}^{1/2}\|D^{2}u\|_{L^{2}}^{1/2}
\end{align*}
对一切 $u\in C^{\infty}_{c}(U)$ 成立。设 $U$ 有界，$\partial U$ 光滑，并证明若 $u\in H^{2}(U)\cap H^{1}_{0}(U)$，此不等式也成立。（提示：取序列 $\{v_{k}\}_{k=1}^{\infty}\subset C^{\infty}_{c}(U)$ 在 $H^{1}_{0}(U)$ 中收敛于 $u$，以及 $\{w_{k}\}_{k=1}^{\infty}\subset C^{\infty}(\overline{U})$ 在 $H^{2}(U)$ 中收敛于 $u$。）

\item \leavevmode
\begin{enumerate}[(a)]
\item 分部积分以证明对 $2\leqslant p<\infty$ 及一切 $u\in C^{\infty}_{c}(U)$ 有
\begin{align*}
\|Du\|_{L^{p}}\leqslant C\|u\|_{L^{p}}^{1/2}\|D^{2}u\|_{L^{p}}^{1/2}.
\end{align*}
（提示：$\int_{U}|Du|^{p}\,\mathrm{d}x=\sum_{i=1}^{n}\int_{U}u_{x_{i}}u_{x_{i}}|Du|^{p-2}\,\mathrm{d}x$。）
\item 证明对 $1\leqslant p<\infty$ 及一切 $u\in C^{\infty}_{c}(U)$ 有
\begin{align*}
\|Du\|_{L^{2p}}\leqslant C\|u\|_{L^{\infty}}^{1/2}\|D^{2}u\|_{L^{p}}^{1/2}.
\end{align*}
\end{enumerate}

\item 举例说明：若对一切 $0<|h|<\frac{1}{2}\operatorname{dist}(V,\partial U)$ 有 $\|D^{h}u\|_{L^{1}(V)}\leqslant C$，并不一定推出 $u\in W^{1,1}(V)$。

\item 给出一个开集 $U\subset\mathbb{R}^{n}$ 及函数 $u\in W^{1,\infty}(U)$ 的例子，使 $u$ 在 $U$ 上不是 Lipschitz 连续的。（提示：取 $U$ 为 $\mathbb{R}^{2}$ 中去掉一条割线的开单位圆盘。）

\item\label{exercise:ch05-18} 验证若 $n>1$，无界函数 $u=\log\log\left(1+\dfrac{1}{|x|}\right)$ 属于 $W^{1,n}(U)$，其中 $U=B^{0}(0,1)$。

\item 固定 $\alpha>0$，令 $U=B^{0}(0,1)$。证明存在常数 $C$，只依赖于 $n$ 与 $\alpha$，使
\begin{align*}
\int_{U}u^{2}\,\mathrm{d}x\leqslant C\int_{U}|Du|^{2}\,\mathrm{d}x,
\end{align*}
只要
\begin{align*}
|\{x\in U\mid u(x)=0\}|\geqslant\alpha,\qquad u\in H^{1}(U).
\end{align*}

\item （Hardy 不等式的变形）证明对每个 $n\geqslant 3$，存在常数 $C$，使对一切 $u\in H^{1}(\mathbb{R}^{n})$ 有
\begin{align*}
\int_{\mathbb{R}^{n}}\frac{u^{2}}{|x|^{2}}\,\mathrm{d}x\leqslant C\int_{\mathbb{R}^{n}}|Du|^{2}\,\mathrm{d}x.
\end{align*}
（提示：对每个 $\lambda\in\mathbb{R}$，$|Du+\lambda\dfrac{x}{|x|^{2}}u|^{2}\geqslant 0$。）

\item 设 $1\leqslant p\leqslant\infty$，$U$ 有界。
\begin{enumerate}[(a)]
\item 证明若 $u\in W^{1,p}(U)$，则 $|u|\in W^{1,p}(U)$。
\item 证明 $u\in W^{1,p}(U)$ 蕴含 $u^{+},u^{-}\in W^{1,p}(U)$，且
\begin{align*}
Du^{+}=
\begin{cases}
Du, & \text{a.e. 于}\ \{u>0\},\\
0, & \text{a.e. 于}\ \{u\leqslant 0\},
\end{cases}
\qquad
Du^{-}=
\begin{cases}
0, & \text{a.e. 于}\ \{u\geqslant 0\},\\
-Du, & \text{a.e. 于}\ \{u<0\}.
\end{cases}
\end{align*}
（提示：$u^{+}=\lim_{\varepsilon\to 0}F_{\varepsilon}(u)$，其中
\begin{align*}
F_{\varepsilon}(z)\triangleq
\begin{cases}
(z^{2}+\varepsilon^{2})^{1/2}-\varepsilon, & z\geqslant 0,\\
0, & z<0.
\end{cases}
\end{align*}
）
\item 证明若 $u\in W^{1,p}(U)$，则 $Du=0$ a.e. 于集合 $\{u=0\}$ 上。
\end{enumerate}

\item 为下面关于“若 $u\in H^{1}(U)$，则 $Du=0$ a.e. 于集合 $\{u=0\}$ 上”的另一证明补充细节。

设 $\phi$ 为光滑、有界且非降函数，使 $\phi'$ 有界且当 $|z|\leqslant 1$ 时 $\phi(z)=z$。令
\begin{align*}
u_{\varepsilon}(x)\triangleq\varepsilon\,\phi(u/\varepsilon).
\end{align*}
证明 $u_{\varepsilon}\rightharpoonup 0$ 弱于 $H^{1}(U)$，因而
\begin{align*}
\int_{U}Du_{\varepsilon}\cdot Du\,\mathrm{d}x=\int_{U}\phi'(u/\varepsilon)|Du|^{2}\,\mathrm{d}x\to 0.
\end{align*}
利用此观察完成证明。

\item 用Fourier变换证明：若 $u\in H^{s}(\mathbb{R}^{n})$ 且 $s>n/2$，则 $u\in L^{\infty}(\mathbb{R}^{n})$，且有界
\begin{align*}
\|u\|_{L^{\infty}(\mathbb{R}^{n})}\leqslant C\|u\|_{H^{s}(\mathbb{R}^{n})},
\end{align*}
其中常数 $C$ 只依赖于 $s$ 与 $n$。

\item 证明若 $u,v\in H^{s}(\mathbb{R}^{n})$ 且 $s>n/2$，则 $uv\in H^{s}(\mathbb{R}^{n})$，且
\begin{align*}
\|uv\|_{H^{s}(\mathbb{R}^{n})}\leqslant C\|u\|_{H^{s}(\mathbb{R}^{n})}\|v\|_{H^{s}(\mathbb{R}^{n})},
\end{align*}
其中常数 $C$ 只依赖于 $s$ 与 $n$。
\end{enumerate}

\end{document}
